\section{Matched-cost execution of lossless neural search}\label{sec:study58}
The preceding experiment asks whether exact fitted neural witnesses alter the returned economic policy. The present experiment holds those witnesses fixed and asks what it costs to compute them. It is a matched-work comparison of two implementations of the same NBO return procedure, not another policy-cost ranking obtained by relabeling identical actors.

\subsection{A prospective, isolated intervention}
The economic laws, discount $15/16$, investment price one, initial uniform law, state-dependent capacities and action spacing remain unchanged. The tasks are $(d,T)=(2,2),(4,4),(8,6)$. At each of the two targets, $q=9/10$ and $q=4/5$ of installed expected cost, each solver starts a separate process from the primitives. Both independently fit the same own-future ReLU continuation, construct the same non-tensor partition, retain the same common actions and apply the same whole-cell verifier. The algebraic solver remains common at the terminal date. Only nonterminal trained-neural action search changes.

Two resource attempts use $(256,32,4)$ and $(1024,64,8)$ training rows, acquired leaves and innovation bins. A complete attempted policy is fixed before its inference begins. The evidence-based return rule and the cumulative looks of 4096, 16384 and 65,536 paths are exactly those in Section~\ref{sec:study57}. A nonpositive upper endpoint for $J(\pi)-qJ(0)$ returns immediately; a positive lower endpoint advances the resource attempt; an unresolved interval advances the look. Exhausting the declared attempts returns a safe policy with an unattained target, not a proof that the target is impossible. Neither attempt is deployed before this return, and each verification uses the same smooth installed reference.

The protocol specifies three fixed-work repetitions, giving 36 complete processes. Both algorithms and all repetitions use one declared training seed and matching inference streams within each task, target and attempt. Exact parameter identities test that the fitted objects match; identical seeds alone are not treated as proof. This design isolates elapsed-time variation at fixed work. It complements, rather than replaces, the preceding three-seed study of different fitted objects.

\input{tables/returns58}

\subsection{Identity is stronger than an unresolved contrast}
The replay reconstructs both action solvers from each distinct saved fitted object, compares their exact objective values and tie choices, and regenerates the whole-cell certificates. It also reconstructs each distinct inference stream and checks the original empirical moments in rational arithmetic. All repeated records are checked against those reconstructions. Table~\ref{tab:returns58} reports the complete target outcomes and actual cost intervals; the full service registry includes every attempted rung and inference look.

Identical partitions and deployed action arrays establish equality of policy costs under the original law, without a confidence interval for their difference. The stronger internal identities also establish the stopping-path conclusion of Corollary~\ref{cor:return-identity58}. Thus a change in work in this comparison cannot be attributed to a different fit, a weaker accuracy target, a changed actor, omitted verification, or fewer declared evidence steps.

The new inference namespace is separate from the earlier study. Within this experiment, repeated use of the same stream does not create additional independent observations. Even counting every repeated interval in the union account, the maximum number is
\[
 36\cdot2\cdot3\cdot3=648<2048.
\]
The inherited logarithm bound 14 therefore gives the declared family error $1/100$. This is a finite, prespecified look family, not an assertion of validity at arbitrary future stopping times. The reported initial-law target is also distinct from the all-state Bellman loss bound, which retains the reference-continuation term and is unchanged between matched solvers.

\subsection{Work, arithmetic, and the economic comparison}
The complete process clock includes imports, source verification, all fitting and search, whole-cell verification, own inference, serialized records, the flushed process log and termination. The timer is external to the measured service. Writing the controller's own clock record is additional orchestration overhead, not a free operation inserted inside a self-timed endpoint. Processes are sequential on a single available core within each task; solver order alternates by repetition and target. Processor frequency is not controlled, and different task runners do not constitute a controlled hardware-scaling experiment.

\input{tables/work58}
\input{tables/arithmetic58}
\input{sections/outcomes58}

The arithmetic ledger concerns nonterminal neural queries only. Terminal search is identical in both modes. Active-region reductions and retained-point counts make the source of any change in rational work observable; all fallback cases are retained. The full vector screen still performs work proportional to the original lattice size. Its fixed-precision operation count is not interchangeable with a count of bit-dependent rational comparisons or cubic-root isolations.

The complete-clock table, rather than an isolated solver timer, determines the observed work comparison at each matched target. Component search times and all three elapsed times are supplied in the supplement. Their medians describe these fixed workloads; three observations do not establish a population timing distribution or a guarantee for another processor.

For an identical returned policy, the price calculation is exactly \eqref{eq:net-work58}. Any supported switch is a resource improvement at unchanged economic output and unchanged certification, not an additional gain in the policy itself. The original common-only, quadratic and tree comparisons remain in Section~\ref{sec:study57}; speeding up one neural subroutine does not establish that the complete neural method is cheaper than every conventional alternative. Likewise, the normalized computation price in the adoption calculation is not an empirical fee calibration.
